\section{评估与奖励：如何把开放任务变成可训练目标}

\subsection{执行式评估（Execution-based Evaluation）}
课程采用脚本化执行评估，而非人工主观打分。任务完成后，由验证脚本检查文件、页面状态、输出内容或系统变量，给出完成度判断。

\begin{importantbox}{开放任务也可以做可计算评估}
即使任务有多条路径，只要结果状态可验证，就能构建自动评测。关键是把 ``任务意图'' 映射为 ``可检查条件''，例如文件存在性、字段值、页面元素状态、命令返回值等。
\end{importantbox}

\begin{figure}[H]
\centering
\includegraphics[width=0.88\textwidth]{frame-04-reward.jpg}
\caption{视频约 00:19:35：讲者介绍执行式奖励与部分完成评分逻辑}
\end{figure}

\subsection{部分完成奖励与难度分层}
课程强调不应只用二元成败。对复杂任务，部分完成分数能更细粒度反映进步，并显著改善训练信号稀疏问题。实践中应配套难度分层，避免模型长期停在低阶子任务。

\begin{knowledgebox}{奖励函数设计建议}
\begin{itemize}
    \item 将目标拆成关键里程碑，设置分段奖励；
    \item 对破坏性动作增加惩罚，约束风险行为；
    \item 对无效重复操作做退火处理，抑制循环轨迹；
    \item 评估脚本版本化管理，保证可复现性。
\end{itemize}
\end{knowledgebox}

\subsection{评测体系的鲁棒性要求}
评测脚本本身也可能被 ``攻击'' 或绕过。课程建议在脚本层加入异常分支、边界检查和反作弊约束，避免模型学习到与任务目标无关的投机策略。

\begin{warningbox}{Reward Hacking 是真实风险}
如果奖励只覆盖局部状态，Agent 可能学会 ``骗分'' 而非 ``完成任务''。例如通过界面缓存制造假完成状态，或触发脚本漏洞规避关键步骤。评估设计必须先于模型迭代。
\end{warningbox}

\subsection{本章小结}
Agent 训练的上限很大程度上由评估系统决定。执行式评估、部分奖励和脚本鲁棒性共同决定了模型能否在开放任务上稳定进步。

